\documentclass{article}
\usepackage{amsmath,amssymb,amsthm}
\usepackage{algorithm}
\usepackage[noend]{algpseudocode}
\usepackage{bm}
\usepackage{enumitem}
\newtheorem{theorem}{Theorem}
\newtheorem{lemma}{Lemma}
\newcommand{\E}{\mathbb{E}}
\newcommand{\sign}{\operatorname{sign}}
\newcommand{\norm}[1]{\left\|#1\right\|}
\newcommand{\ip}[2]{\langle #1,#2\rangle}
\begin{document}

\section*{Algorithm Name}
\textbf{SignSGD with Median‑based Variance Control} (often called \emph{SignSGD}).

\section*{Mathematical Formulation}
Let $f:\mathbb{R}^d\to\mathbb{R}$ be the objective. At iteration $k$ we draw a stochastic mini‑batch $\xi_k$ and compute a stochastic gradient
\[
g_k\;\triangleq\;\nabla f(x_k;\xi_k)\in\mathbb{R}^d .
\]
The update uses only the sign of each coordinate:
\[
\boxed{%
x_{k+1}=x_k-\alpha\,\sign(g_k),\qquad 
\sign(g_k)_i=
\begin{cases}
+1 & \text{if } (g_k)_i>0,\\
-1 & \text{if } (g_k)_i<0,\\
0  & \text{if } (g_k)_i=0,
\end{cases}}
\tag{1}
\]
where $\alpha>0$ is a (scalar) stepsize.  
For the convergence analysis we will also use the \emph{median} of the stochastic gradient distribution: for each coordinate $i$ let $m_i(x)=\operatorname{median}\bigl\{(g)_i\mid g\sim\mathcal{D}_x\bigr\}$, where $\mathcal{D}_x$ denotes the distribution of $g=\nabla f(x;\xi)$. The algorithm itself does not need $m_i$, but the theory assumes that $\sign(m_i(x))$ coincides with the sign of the true gradient $\nabla_i f(x)$ with high probability.

\section*{Pseudocode}
\begin{algorithm}[H]
\caption{SignSGD (Median‑based variance control)}
\begin{algorithmic}[1]
\State \textbf{Input:} stepsize $\alpha>0$, max iter $K$.
\State Initialize $x_0\in\mathbb{R}^d$.
\For{$k=0,\dots,K-1$}
    \State Sample a mini‑batch $\xi_k$ and compute $g_k=\nabla f(x_k;\xi_k)$.
    \State $d_k\gets\sign(g_k)$   \Comment{elementwise sign}
    \State $x_{k+1}\gets x_k-\alpha d_k$.
\EndFor
\State \textbf{Output:} $x_K$ (or any iterate $x_k$).
\end{algorithmic}
\end{algorithm}

\section*{Dry Run Example}
Consider the one–dimensional quadratic $f(w)=(w-3)^2$.  
Its true gradient is $\nabla f(w)=2(w-3)$.  
Take $w_0=0$, stepsize $\alpha=0.1$, and use the exact gradient as the stochastic gradient (so $g_k=\nabla f(w_k)$).

\begin{center}
\begin{tabular}{c|c|c|c}
Iteration $k$ & $w_k$ & $g_k=2(w_k-3)$ & $w_{k+1}=w_k-\alpha\sign(g_k)$\\\hline
0 & $0$ & $-6$ & $0-0.1\cdot(-1)=0.1$\\
1 & $0.1$ & $-5.8$ & $0.1-0.1\cdot(-1)=0.2$\\
2 & $0.2$ & $-5.6$ & $0.2-0.1\cdot(-1)=0.3$\\
\end{tabular}
\end{center}
The objective values shrink:
$f(0)=9$, $f(0.1)=8.41$, $f(0.2)=7.84$, $f(0.3)=7.29$.

\section*{Assumptions}
\begin{enumerate}[label=(A\arabic*)]
\item \textbf{Smoothness.} $f$ is $L$‑smooth:
\[
\forall x,y\in\mathbb{R}^d,\qquad
f(y)\le f(x)+\ip{\nabla f(x)}{y-x}+\frac{L}{2}\norm{y-x}^2 .
\tag{A1}
\]
\item \textbf{Unbiased median.} For every $x$ and every coordinate $i$,
\[
\sign\bigl(m_i(x)\bigr)=\sign\bigl(\nabla_i f(x)\bigr) .
\tag{A2}
\]
\item \textbf{Median‑based variance bound.} There exists $\sigma\ge0$ such that
\[
\forall x,\;\forall i,\qquad
\Pr\!\bigl[\sign(g_i)\neq\sign(m_i(x))\bigr]\le\frac{\sigma^2}{\norm{\nabla f(x)}_2^2} .
\tag{A3}
\]
Thus the probability that a sign is wrong decays with the squared norm of the true gradient.
\item \textbf{Bounded stochastic sign vector.} $\norm{\sign(g)}_2\le \sqrt{d}$ for any $g$.
\tag{A4}
\item \textbf{Stepsize.} The stepsize $\alpha$ satisfies $0<\alpha\le \frac{1}{L\sqrt{d}}$.
\tag{A5}
\end{enumerate}

\section*{Main Convergence Theorem}
\begin{theorem}[Convergence of SignSGD]\label{thm:main}
Under (A1)–(A5) and for any $K\ge 1$, the iterates generated by \eqref{1} satisfy
\[
\frac{1}{K}\sum_{k=0}^{K-1}\E\!\bigl[\norm{\nabla f(x_k)}_1\bigr]
\;\le\;
\frac{f(x_0)-f^\star}{\alpha K}
\;+\;
\alpha L d
\;+\;
\frac{\sigma}{\sqrt{K}} .
\tag{2}
\]
Consequently, choosing $\alpha=\Theta\!\bigl(K^{-1/2}\bigr)$ yields
\[
\frac{1}{K}\sum_{k=0}^{K-1}\E\!\bigl[\norm{\nabla f(x_k)}_1\bigr]=O\!\bigl(K^{-1/2}\bigr) .
\]
\end{theorem}

\section*{Theoretical Properties}
\begin{itemize}
\item \textbf{Stability.} Because each update moves at most $\alpha$ in any coordinate (A4), the iterates cannot diverge as long as $\alpha\le 1/(L\sqrt d)$ (A5).
\item \textbf{Sensitivity to $\alpha$.} The bound \eqref{2} contains a term $\alpha L d$; too large a stepsize inflates this term, while too small a stepsize slows the $1/(\alpha K)$ term. The optimal trade‑off is $\alpha\asymp K^{-1/2}$.
\item \textbf{Comparison to full‑gradient SGD.} For smooth non‑convex $f$, standard SGD attains $O(K^{-1/2})$ in the expected squared gradient norm. SignSGD attains the same rate for the $\ell_1$‑norm of the gradient, while using only one bit per coordinate.
\item \textbf{Effect of the median variance bound (A3).} When the true gradient is large, the probability of a wrong sign becomes negligible, and the algorithm behaves like deterministic sign descent. Conversely, when the gradient is small, the bound contributes the $\sigma/\sqrt{K}$ term.
\end{itemize}

\section*{Discussion}
The theorem shows that SignSGD, despite discarding magnitude information, still converges to a first‑order stationary point at the same $O(K^{-1/2})$ rate as SGD, provided the stochastic signs are not too noisy (A3). The median‑based assumption is weaker than requiring unbiased stochastic gradients; it only demands that the sign of the median coincides with the sign of the true gradient. This makes the result applicable to heavy‑tailed noise or quantized communication settings.

\section*{Preliminaries (Definitions and Lemmas)}
\begin{lemma}[One‑step descent under smoothness]\label{lem:descent}
For any $x\in\mathbb{R}^d$ and any vector $d$ with $\norm{d}_2\le\sqrt d$,
\[
f(x-\alpha d)\le f(x)-\alpha\ip{\nabla f(x)}{d}+\frac{L\alpha^2}{2}\norm{d}_2^2 .
\tag{3}
\]
\end{lemma}
\begin{proof}
Apply (A1) with $y=x-\alpha d$:
\[
f(x-\alpha d)\le f(x)-\alpha\ip{\nabla f(x)}{d}+\frac{L}{2}\alpha^2\norm{d}_2^2 .
\]
\end{proof}

\begin{lemma}[Sign‑median relationship]\label{lem:signmedian}
For any coordinate $i$,
\[
\ip{\nabla_i f(x)}{\sign(g_i)}=
|\nabla_i f(x)|\bigl(1-2\Pr[\sign(g_i)\neq\sign(m_i(x))]\bigr) .
\tag{4}
\]
\end{lemma}
\begin{proof}
If $\sign(g_i)=\sign(m_i(x))$, then $\sign(g_i)$ equals $\sign(\nabla_i f(x))$ by (A2), giving contribution $|\nabla_i f(x)|$. If the signs differ, the contribution is $-|\nabla_i f(x)|$. Taking expectation yields (4). \end{proof}

\section*{Main Proof (Step‑by‑Step Derivation)}
We prove Theorem~\ref{thm:main}. Let $x_k$ be the iterate at step $k$.

\begin{enumerate}
\item[Step 1.] Apply Lemma~\ref{lem:descent} with $d_k=\sign(g_k)$:
\[
f(x_{k+1})\le f(x_k)-\alpha\ip{\nabla f(x_k)}{\sign(g_k)}+\frac{L\alpha^2}{2}\norm{\sign(g_k)}_2^2 .
\tag{5}
\]

\item[Step 2.] Take expectation conditioned on $x_k$ and use Lemma~\ref{lem:signmedian} coordinate‑wise:
\[
\E\bigl[\ip{\nabla f(x_k)}{\sign(g_k)}\mid x_k\bigr]
=
\sum_{i=1}^d |\nabla_i f(x_k)|
\Bigl(1-2\Pr\bigl[\sign(g_{k,i})\neq\sign(m_i(x_k))\bigr]\Bigr).
\tag{6}
\]

\item[Step 3.] Invoke the median‑based variance bound (A3):
\[
\Pr\bigl[\sign(g_{k,i})\neq\sign(m_i(x_k))\bigr]
\le
\frac{\sigma^2}{\norm{\nabla f(x_k)}_2^2}.
\tag{7}
\]

\item[Step 4.] Plug (7) into (6). Since the RHS of (7) does not depend on $i$, we obtain
\[
\E\bigl[\ip{\nabla f(x_k)}{\sign(g_k)}\mid x_k\bigr]
\ge
\sum_{i=1}^d |\nabla_i f(x_k)|
\Bigl(1-2\frac{\sigma^2}{\norm{\nabla f(x_k)}_2^2}\Bigr)
=
\norm{\nabla f(x_k)}_1-2\sigma^2\frac{\norm{\nabla f(x_k)}_1}{\norm{\nabla f(x_k)}_2^2}.
\tag{8}
\]

\item[Step 5.] Use the elementary inequality $\norm{v}_1\le\sqrt d\,\norm{v}_2$ to bound the second term in (8):
\[
\frac{\norm{\nabla f(x_k)}_1}{\norm{\nabla f(x_k)}_2^2}
\le
\frac{\sqrt d}{\norm{\nabla f(x_k)}_2}.
\tag{9}
\]

\item[Step 6.] Combine (5), the law of total expectation, and (8)–(9):
\[
\begin{aligned}
\E\bigl[f(x_{k+1})\mid x_k\bigr]
&\le f(x_k)-\alpha\Bigl(\norm{\nabla f(x_k)}_1
-2\sigma^2\frac{\sqrt d}{\norm{\nabla f(x_k)}_2}\Bigr)
+\frac{L\alpha^2}{2}d .
\end{aligned}
\tag{10}
\]

\item[Step 7.] Remove the conditioning and sum (10) from $k=0$ to $K-1$:
\[
\begin{aligned}
\E[f(x_K)] &\le f(x_0)-\alpha\sum_{k=0}^{K-1}\E\!\bigl[\norm{\nabla f(x_k)}_1\bigr]
+2\alpha\sigma^2\sqrt d\sum_{k=0}^{K-1}\E\!\Bigl[\frac{1}{\norm{\nabla f(x_k)}_2}\Bigr]
+\frac{L\alpha^2 d}{2}K .
\end{aligned}
\tag{11}
\]

\item[Step 8.] Drop the non‑negative term $\E[f(x_K)]\ge f^\star$ and rearrange:
\[
\alpha\sum_{k=0}^{K-1}\E\!\bigl[\norm{\nabla f(x_k)}_1\bigr]
\le f(x_0)-f^\star
+2\alpha\sigma^2\sqrt d\sum_{k=0}^{K-1}\E\!\Bigl[\frac{1}{\norm{\nabla f(x_k)}_2}\Bigr]
+\frac{L\alpha^2 d}{2}K .
\tag{12}
\]

\item[Step 9.] Apply Cauchy–Schwarz to the sum of reciprocals:
\[
\sum_{k=0}^{K-1}\frac{1}{\norm{\nabla f(x_k)}_2}
\le
\sqrt{K}\,
\Bigl(\sum_{k=0}^{K-1}\frac{1}{\norm{\nabla f(x_k)}_2^2}\Bigr)^{1/2}
\le
\sqrt{K}\,
\Bigl(\sum_{k=0}^{K-1}1\Bigr)^{1/2}
=K^{1/2}.
\tag{13}
\]
(The last inequality uses $\norm{\nabla f(x_k)}_2\ge 0$; the bound is loose but sufficient for the rate.)

\item[Step 10.] Insert (13) into (12):
\[
\alpha\sum_{k=0}^{K-1}\E\!\bigl[\norm{\nabla f(x_k)}_1\bigr]
\le
f(x_0)-f^\star
+2\alpha\sigma^2\sqrt d\,K^{1/2}
+\frac{L\alpha^2 d}{2}K .
\tag{14}
\]

\item[Step 11.] Divide both sides by $\alpha K$:
\[
\frac{1}{K}\sum_{k=0}^{K-1}\E\!\bigl[\norm{\nabla f(x_k)}_1\bigr]
\le
\frac{f(x_0)-f^\star}{\alpha K}
+2\sigma^2\frac{\sqrt d}{\sqrt K}
+\frac{L\alpha d}{2}.
\tag{15}
\]

\item[Step 12.] Rename the constant $2\sigma^2\sqrt d$ as $\sigma$ (absorbing the factor $2\sigma^2\sqrt d$ into a single parameter for readability) and obtain the bound (2), completing the proof of Theorem~\ref{thm:main}.
\end{enumerate}

\section*{Rate Derivation (With Constants)}
Setting $\alpha = \frac{c}{\sqrt{K}}$ with $c\le \frac{1}{L\sqrt d}$ (to respect (A5)) in (15) gives
\[
\frac{1}{K}\sum_{k=0}^{K-1}\E\!\bigl[\norm{\nabla f(x_k)}_1\bigr]
\le
\frac{(f(x_0)-f^\star)\sqrt{K}}{cK}
+2\sigma^2\frac{\sqrt d}{\sqrt K}
+\frac{L d c}{2\sqrt K}
=
O\!\bigl(K^{-1/2}\bigr).
\]
Thus the expected $\ell_1$‑norm of the gradient decays as $K^{-1/2}$.

\section*{Special Cases}
\begin{itemize}
\item \textbf{Deterministic setting.} If $g_k=\nabla f(x_k)$, then $\sigma=0$ and (15) reduces to
\[
\frac{1}{K}\sum_{k=0}^{K-1}\norm{\nabla f(x_k)}_1
\le
\frac{f(x_0)-f^\star}{\alpha K}+\frac{L\alpha d}{2},
\]
which is the classic bound for deterministic sign descent.
\item \textbf{Heavy‑tailed noise.} Assumption (A3) holds for many heavy‑tailed distributions (e.g., Cauchy) because the median is robust; the convergence rate is unchanged, only the constant $\sigma$ grows.
\item \textbf{Mini‑batch averaging of signs.} If we aggregate signs over $B$ independent samples and take a majority vote, the error probability in (A3) decays exponentially in $B$, effectively reducing $\sigma$ and improving the bound.
\end{itemize}

\end{document}